\slajdid{08-s065}
\begin{frame}[label=08-s065]
\gusttekst
\[p(t)=p(\infty)+\bigl(p(t_0^+)-p(\infty)\bigr)e^{-\frac{t-t_0}{\tau}}\qquad za\ t\geq t_0\]
\centering 1. $\tau$\par\raggedright
Određujemo kao \textbf{proizvod L i provodnosti} koju vidi ta induktivnost tako što sve nezavisne naponske izvore kratko spojimo, nezavisne strujne izvore uklonimo iz kola, ostavimo otvorene veze, pri čemu zavisni izvori ostaju.\par\medskip
\centering 2. $p(t_0^+)$\par\raggedright
Određujemo na osnovu saznanja da je u pitanju realno kolo i da struja kroz induktivnost ne može trenutno da se promeni.
\[i_L(t_0^+)=i_L(t_0^-)\]
\centering 3. $p(\infty)$\par\raggedright
Određujemo na osnovu saznanja da su u beskonačnosti završeni prelazni procesi odnosno da nema promene stuje kroz induktivnost. A uslov da nema promene struje kroz induktivnost je
\[u_L(\infty)=0\]
\end{frame}
